\section{Two-plane carriers and ordinary orthogonal groups in odd characteristic}
\label{sec:orthogonal-even-odd}

\paragraph{Status and boundary.}
This is a separate new author derivation.  The full written proof
and exact supplement pass the author check; independent review
is pending and the result is not promoted.  It concerns
ordinary representations of the natural group $\Omega(V)$.  It makes
no assertion for representations of a spin cover which do not factor
through this natural group.  Earlier canonical sections are unchanged.

We use the exact constants of D-OL-14:
\[
 \beta=\frac16-\frac1{10000}=\frac{4997}{30000},\qquad
 \delta=\frac{1-(19/20)^9}{2},\qquad
 \frac16<\delta<\frac3{16}<\frac7{32}.
 \tag{OE.1}
\]
For a complex operator of degree $a$, write
$\nu(T)=a-\max_\lambda\dim\ker(T-\lambda I)$.

\subsection*{D-OL-16. A two-plane orthogonal carrier}

\begin{lemma}[D-OL-16]
Let $F=\mathbf F_q$, with $q$ odd, and let $V$ be a nondegenerate
quadratic space of even dimension $d\ge24$.  Put $S=\Omega(V)$.
Suppose that a linear similitude $H$ satisfies
\[
       H^2=cI,\qquad Q(Hv)=\mu Q(v),\qquad
       c\notin(F^*)^2.
       \tag{OE.2}
\]
Let $T$ be a finite-projective-order compatible normalizer of a
nontrivial ordinary irreducible complex $S$-module of degree $a$,
implementing conjugation by $H$.  Then
\[
                              \nu(T)\ge7a/32.
                              \tag{OE.3}
\]
No extension of this module to a larger orthogonal group is used.
\end{lemma}

\begin{proof}
Since $\mu^2=c^2$, write $\varepsilon=\mu/c\in\{1,-1\}$.
Let $E=F[t]/(t^2-c)$, and make $V$ an $E$-space by letting $t$
act as $H$.  Write $B$ for the polar symmetric form, so
$Q(v)=B(v,v)/2$.  The adjoint identity is $H^*=\varepsilon H$.
Define
\[
 b(v,w)=\frac12\bigl(B(v,w)+t^{-1}B(Hv,w)\bigr).
 \tag{OE.4}
\]
If $\varepsilon=1$, this is a nondegenerate symmetric $E$-bilinear
form.  If $\varepsilon=-1$, it is a nondegenerate $E/F$-Hermitian
form, linear in its first variable.  Indeed multiplication by $t$
in the first variable multiplies (OE.4) by $t$; in the second
variable it multiplies it by $\varepsilon t$.  Symmetry of $B$
gives the indicated symmetry or conjugate symmetry, and
$\operatorname{Tr}_{E/F}b=B$ proves nondegeneracy.

The finite-field symmetric classification and the finite-field
Hermitian hyperbolic-pair theorem supply a hyperbolic pair for $b$:
the $E$-dimension is at least twelve, while the respective
anisotropic dimensions are at most two and one
\cite[pp.~138--140; Lemma~10.2, Corollary~10.3 and the ensuing
decomposition, pp.~116--117]{OS:Taylor1992}.
Its two $E$-lines give opposite, totally singular, $H$-stable
$F$-planes $L,L'$.  Their cross pairing is nondegenerate.  Choose
dual $F$-bases and write
\[
 V=L\oplus W\oplus L',\quad
 Q(x,w,y)=x^{\mathsf t}y+\tfrac12w^{\mathsf t}Bw,
 \quad N=\dim W=d-4\ge20.
 \tag{OE.5}
\]
Here and below $B$ also denotes the nonsingular matrix on $W$.
In these coordinates
\[
 H=\operatorname{diag}(A,h,\mu A^{-\mathsf t}),\quad
 A^2=cI_2,\quad h^2=cI_N,\quad
 h^{\mathsf t}Bh=\mu B,\quad\det A=-c.
 \tag{OE.6}
\]

The unipotent radical fixing $L$ is the following actual subgroup:
\[
 u(X,Z)=
 \begin{pmatrix}
 I_2&X&Z-\tfrac12XB^{-1}X^{\mathsf t}\\
 0&I_N&-B^{-1}X^{\mathsf t}\\
 0&0&I_2
 \end{pmatrix},\qquad
 X\in M_{2,N}(F),\quad Z^{\mathsf t}=-Z.
 \tag{OE.7}
\]
Direct multiplication gives
\[
 u(X,Z)u(Y,R)=u\left(X+Y,
 Z+R+\tfrac12(YB^{-1}X^{\mathsf t}-XB^{-1}Y^{\mathsf t})\right).
 \tag{OE.8}
\]
Every such matrix is an isometry, and this is a group of exponent
$p=\operatorname{char}F$.  Thus its determinant and spinor norm
are trivial and it lies in $S$.  The alternating pairing on the
$2N$-dimensional $X$-space given by the $(1,2)$ entry in (OE.8)
is nondegenerate.  Consequently this group $Q_L$ is the Heisenberg
group of order $q^{2N+1}$, with center $X=0$.

We justify the fixed-space estimate for this larger radical.
Let $e_1,e_2$ be the chosen basis of $L$ and $f_1,f_2$ its dual.
The radicals $Q_L,Q_{L'}$ contain the two full line radicals on
$Fe_2\oplus W\oplus Ff_2$, with parameters ranging over $W$.
D-OL-06 therefore puts
$\Omega(Fe_2\oplus W\oplus Ff_2)$ in their generated subgroup.
It fixes $e_1,f_1$ and is transitive on the nonzero singular
vectors of its natural space.  The center of $Q_L$ contains
the line-radical elements $u_{e_1}(s e_2)$, for every $s\in F$.
Conjugating by that smaller $\Omega$ gives the parameters $s f_2$;
the rows of $X$ already give every parameter in $W$.  Additivity
of the line radical now gives its full parameter space
$W\oplus Fe_2\oplus Ff_2$.  The same argument applies to $f_1$.
These two full opposite line radicals generate $S$, again by
D-OL-06.  Hence
\[
                    \langle Q_L,Q_{L'}\rangle=S.
                    \tag{OE.9}
\]
The radicals are conjugate inside $S$: exchange the two
hyperbolic pairs, and, if necessary, correct the determinant and
spinor norm in their nondegenerate perpendicular space.  That
space represents both nonzero square classes, so its reflections
give both corrections.  An invariant vector for both radicals
would be $S$-invariant.  Irreducibility and nontriviality thus give
\[
                        t_0:=\dim\rho^{Q_L}\le a/2.
                        \tag{OE.10}
\]

Conjugating (OE.7) by (OE.6) gives exactly
\[
             X\longmapsto AXh^{-1},\qquad
             Z\longmapsto\mu^{-1}AZA^{\mathsf t}.
             \tag{OE.11}
\]
The action on the one-dimensional center is the scalar
$\det A/\mu=-c/\mu=-\varepsilon$.
The action on $X$ is an involution.  Its fixed matrices are
$E$-linear maps $W\to L$, and therefore form an $F$-space of
dimension $N$.  Dually, using the perfect trace pairing on
matrices, the dual weight action on $Y\in\operatorname{Hom}_F(L,W)$
is $Y\mapsto hYA^{-1}$.  Its fixed space also has dimension $N$.
Every nonzero fixed $Y$ has rank two, since its kernel is an
$A$-stable $F$-subspace of the irreducible two-dimensional
$F[A]$-module $L$.

We next bound the fixed fraction of every nonzero linear-weight
orbit, not just the fixed fraction of the whole dual space.
The Levi subgroup $\Omega(W)$ lies in $S$ and acts on $Y$ by
left multiplication.  For rank two, its orbits consist precisely
of the ordered independent pairs of vectors of $W$ with a fixed
Gram matrix.  Witt's theorem extends an isometry between their
spans, even when that span is degenerate
\cite[Theorem~7.4 and complete proof, pp.~57--59]{OS:Taylor1992}.
An extension in $O(W)$ can be corrected into $\Omega(W)$ while
fixing the pair: its perpendicular space contains a nondegenerate
subspace of dimension at least $N-4$, whose reflections realize
all determinant/spinor-norm pairs.  This proves the orbit assertion.

Here is a uniform lower bound for each such orbit.  There are
at least
\[
                    q^{N-1}-q^{\lceil N/2\rceil}-1
\]
choices for its nonzero first vector, by the quadratic Gauss
estimate of D-OL-06.  For each choice the prescribed cross
pairing cuts out an affine hyperplane.  The restricted polar
rank there is at least $N-2$.  Diagonalizing the quadratic
polynomial on this affine hyperplane gives the same Gauss bound
also in the presence of linear terms: a nonzero linear term on
the radical makes the sum zero, and otherwise the absolute sum
is bounded by the usual rank estimate.  Thus the number of
second vectors of the prescribed norm and cross pairing,
after removing at most $q$ dependent vectors, is at least
\[
                    q^{N-2}-q^{N/2}-q.
\]
For $q\ge3$, $N\ge20$, each displayed lower bound exceeds
half its leading term.  Every rank-two orbit therefore has
more than $q^{2N-3}/4$ vectors.  Its fixed fraction under
(OE.11) is at most
\[
                        4q^{-N+3}<1/8.
                        \tag{OE.12}
\]
Rank-one weights have no fixed vector at all.  Weight
multiplicities are constant on these Levi orbits.  Hence on
the sum of nontrivial linear-character spaces of $Q_L$, every
eigenspace of $T$ has codimension at least $7/16$ of that
sum's dimension: the nonfixed weights are paired, and a
paired block contributes at most half its dimension to one
eigenvalue.  This remains true with arbitrary scalar phases
on the weight spaces.

For clarity, the nonlinear Heisenberg characters are also
treated with their multiplicities.  The explicit construction
and matrix-algebra proof in D-HOS-01 show that, for every
nontrivial additive central character $\psi$, there is a unique
irreducible carrier $U_\psi$ of dimension $q^N$, and its Weyl
operators $L_X$ form a basis of $\operatorname{End}(U_\psi)$.
If $\varepsilon=1$, (OE.11) sends $\psi$ to $\psi^{-1}$.
They are distinct because the characteristic is odd.  The
corresponding isotypic spaces are paired, again giving support
at least one half of their combined dimension.

If $\varepsilon=-1$, the central character is fixed.  Let
$\mathcal A:X\mapsto AXh^{-1}$ be the symplectic involution
on the $X$-space.  Its fixed space has dimension $N$.  The
carrier intertwiner $\Xi$, normalized to satisfy $\Xi^2=I$,
obeys the exact identity
\[
 \Xi L_X\Xi^{-1}=L_{\mathcal A X},\qquad
 |\operatorname{Tr}\Xi|^2
       =\operatorname{Tr}(\operatorname{Ad}\Xi)
       =q^N.
 \tag{OE.13}
\]
The first equality has no extra phase, since (OE.11) preserves
the actual zero-central-coordinate section in (OE.7)--(OE.8).
Thus the support fraction of $\Xi$ is at least
$(1-q^{-N/2})/2>7/16$.  On an isotypic space the operator is
$\Xi\otimes C$ by the matrix-algebra commutant.  Eigenvalue
convolution preserves this support fraction; diagonalizability
follows by rescaling the given finite-projective-order $T$ to
finite order.  No multiplicity operator has been discarded.

The trivial $Q_L$-isotypic space has dimension $t_0$.  Combining
all the preceding summands and (OE.10) proves
$\nu(T)\ge(7/16)(a-t_0)\ge7a/32$, as required.
\end{proof}

\subsection*{D-OL-17. All compatible normalizers and the exceptional reflections}

\begin{theorem}[D-OL-17]
Let $q$ be odd, let $d\ge1136$ be even, and let $\rho$ be any
nontrivial ordinary irreducible complex representation of
$S=\Omega_d^\pm(q)$, of degree $a$.  If $q\ge5$, every nonscalar
compatible finite-projective-order normalizer satisfies
$\nu(T)>\delta a$.  If $q=3$, every such normalizer satisfies
$\nu(T)\ge\beta a$, and
\[
 \nu(T)\le\delta a\quad\Longrightarrow\quad
 \text{its natural implementing projective map is a reflection}.
 \tag{OE.14}
\]
In the last case $T$ itself has projective order two.
\end{theorem}

\begin{proof}
The complete natural semilinear-normalizer description,
including represented central quotients, is D-OL-07.  Thus a
normalizer is implemented by a semilinear similitude $H$.
First suppose that $T$ has prime projective order $l$.  The
centralizer and perfectness argument in D-OL-07 gives $H^l=cI$.
If its field part is nontrivial, the full mixed norm/form descent
of (ON.7)--(ON.9) gives support at least $3a/16>\delta a$.
This includes diagonal--field and graph--field components.
We may therefore suppose that $H$ is linear, with multiplier $\mu$.

If $l$ is odd, then $\mu^l=c^2$.  The scalar
\[
                    z=c/\mu^{(l-1)/2}
                    \quad\text{satisfies}\quad
                    z^2=\mu,\quad z^l=c.
                    \tag{OE.15}
\]
Consequently $g=z^{-1}H$ is an isometry of order $l$.
Its determinant and spinor norm, being of orders dividing both
$l$ and two, are trivial.  Thus $g\in S$, and Schur's lemma
makes $T$ a scalar times $\rho(g)$.

If $g$ fixes a hyperbolic pair, use its complementary line
radical as in (OO.3)--(OO.7).  Its natural moving rank there
is at least two, because an odd-characteristic rank-one
isometry is a reflection and has determinant minus one.
The all-rank Gauss estimate therefore gives a fixed-weight
fraction less than $1/8$, and support at least $7a/32$.
If no hyperbolic pair is fixed, a semisimple fixed space has
dimension at most two.  For a characteristic-order unipotent,
its fixed space has polar rank at least $d-2r$, where
$r=\operatorname{rank}(g-I)$; absence of a hyperbolic pair
therefore gives $r\ge(d-2)/2\ge567$.  The negative eigenspace
is zero in both cases.  In every remaining case the natural
support restricted to $\mu_2=\{1,-1\}$ is at least $567$.
Inflation to the ordinary spin cover and the exact published
LST estimate yield
\[
                    |\chi(g)|/a<(19/20)^9,
                    \qquad\nu(\rho(g))>\delta a.
                    \tag{OE.16}
\]
The source is \cite[Theorem~4.3.6, (4.3.2), pp.~1917--1918;
complete proof pp.~1918--1923; Remark~4.3.7, p.~1923]{HOS:LST2011},
with threshold $567=7\cdot9^2$.  The remark permits the
restricted scalar set $\mu_2$ for spin groups.  Ordinary
inflation is irreducible and no faithful-character assumption
is imposed.  This is the same checked source interface as
(OO.8)--(OO.10), with even $d\ge1136$.

It remains to treat $l=2$.  Now $H^2=cI$ and $\mu=\pm c$.
If $c$ is nonsquare, D-OL-16 applies.  Suppose that $c=z^2$
and replace $H$ by $z^{-1}H$, so $H^2=I$ and $\mu=\pm1$.
If $\mu=1$, the two eigenspaces are nondegenerate.  At least
one contains a hyperbolic pair.  If its eigenvalue is
$\epsilon$, the induced action on the complementary line
radical is $A=\epsilon H|_W$.  It is a nonidentity involution.
If its moving rank is at least two, the preceding $1/8$
bound gives $7a/32$.  If its moving rank is one, the original
projective map is a reflection.  The rank-one Gauss bound is
\[
 \frac{q^{-1}+q^{-9}}{1-q^{-9}-q^{-19}}.
 \tag{OE.17}
\]
For $q\ge5$ it is less than $1/4$, giving $3a/16$.
For $q=3$ it is less than $1-4\beta$, giving $\beta a$.

If instead $\mu=-1$, the two eigenspaces of $H$ are opposite
totally singular spaces of equal dimension.  Choose a
hyperbolic pair with its two vectors in the different
eigenspaces.  The induced action on the line radical is
$A=-H|_W$ after choosing the eigenvalue of the first vector
to be one.  This is an involutory anti-isometry.  After
identifying dual weights using the polar form, the dual
action is $-A$.  Both $A$ and $-A$ have totally singular
fixed spaces of dimension $N/2$, where $N=d-2$.
A nonzero fixed weight can therefore occur only
on the norm-zero Levi orbit.  The fraction fixed in that
orbit is at most
\[
 \frac{q^{N/2}}{q^{N-1}-q^{\lceil N/2\rceil}-1}<\frac18
 \quad(N\ge20,q\ge3).
 \tag{OE.18}
\]
All other norm orbits have no fixed vector.  Opposite line
radicals and paired weights give $7a/32$ once more.  The
sign and scaling in this argument follow directly by
conjugating the line-radical matrix by
$\operatorname{diag}(1,H|_W,-1)$; its parameter action is
$-H|_W$.  Thus no assertion of isometric action was used
for this anti-isometric case.

These arguments prove the asserted bounds for prime
projective order.  Taking a nonscalar prime-order power
never increases projective support, so the full bound
for $q\ge5$ and the lower bound $\beta a$ for $q=3$ follow.
We must prove the stronger all-order classification (OE.14).
For $q=3$ there is no field component.  If the projective
order of $T$ has an odd prime divisor, the corresponding
odd-prime power already has the stronger bound (OE.16).
It remains to consider order $2^s$.  Its involution power
must be a projective reflection if $\nu(T)\le\delta a$.
Suppose that $s\ge2$, and let $H$ be the natural linear lift.
Then
\[
                         H^{2^{s-1}}=\eta r_v
                         \quad(\eta\in F^*).
                         \tag{OE.19}
\]
The map $H$ commutes with this power, hence preserves the
nonsingular one-dimensional $(-1)$-eigenspace $Fv$ of $r_v$.
If $Hv=\lambda v$, its similitude multiplier is $\lambda^2$.
Rescale $H$ by $\lambda^{-1}$ to an isometry.  Its even
power has determinant one.  But any scalar multiple of
$r_v$ which is an isometry is $\pm r_v$, and in even
dimension both have determinant minus one.  This
contradicts (OE.19).  Therefore $s=1$, proving the entire
all-order statement, not merely a statement about one
prime-order power.
\end{proof}

\subsection*{D-OL-18. The actual source obstruction}

\begin{theorem}[D-OL-18]
A nonempty irreducibly acting normal central product of nontrivial
ordinary orthogonal factor modules with odd $q_i$ and even
$d_i\ge1136$ cannot supply period-image rank nine for its
actual six-dimensional multiplicity variation.  Mixed fields,
mixed dimensions, represented central quotients, arbitrary
scalars and all compatible factor permutations are allowed.
No selfduality, determinant or inducing-field equality is
assumed.

In particular, in the main transitive orthogonal-core setting,
every surviving ordinary factor has natural dimension at most
$1134$.  For odd natural dimension the sharper bound is
$1133$.  The separate spin-representation boundary is retained.
\end{theorem}

\begin{proof}
Let $e=\prod_i a_i$.  The degree bound
\cite[Theorem~3.7, pp.~165--167]{OL:GLT2024} gives
$a_i\ge q_i^{d_i-4}(q_i-1)>360$.
Every nonscalar factor-preserving operator has support at
least $\beta e$ by D-OL-17 and eigenvalue convolution.
A nontrivial factor permutation has support at least $e/4$;
a nonidentity scalar has full actual moving rank $e$.

Call an actual inertia good if its moving rank is at least
$\delta e$, and bad otherwise.  A bad inertia fixes every
factor, and every active factor must have $q_i=3$ and be
implemented by a projective reflection.  Its factor
operator has projective order two, by (OE.14).
If two or more factors were active, normalize two of these
operators to involutions.  Their two eigenvalue proportions
lie in $[\beta,1-\beta]$.  The maximum proportion after
convolving them is at most
$(1-\beta)^2+\beta^2$, attained at a corner of this square.
Consequently their support is at least
\[
                       2\beta(1-\beta)e>e/4>\delta e.
                       \tag{OE.20}
\]
Further tensor factors cannot increase the maximum
eigenvalue proportion.  A bad inertia therefore acts
nontrivially on exactly one factor, there by a projective
reflection.  Its global scalar must make one of the two
eigenvalues equal to one, or its actual moving rank is $e$.

The Euler identity for the actual intermediate cover is
$6=2e(h-1)+D$.  Since $e>360$, genus at least two is
impossible.  At genus one $D=6$, so there are at most six
genuine marks and the actual period-image rank is at most
six.  Rank nine therefore requires $h=0$, at least twelve
genuine inertia points, and $D=2e+6$.
The exact inequalities
\[
                  13\beta-2>1/7,
             \qquad11\beta+\delta-2>1/60
             \tag{OE.21}
\]
show first that thirteen marks would imply $e<42$, and
then that twelve marks with one good inertia would imply
$e<360$.  Thus exactly twelve inertias would all have to
be bad.

We now use their actual generation.  They all preserve
the factors.  Fix one normal factor.  The conjugation
image of the inducing group on it contains its entire
inner automorphism group, because the normal factor is
a subgroup of the inducing image.  This conjugation image
is generated by the projections of the twelve actual
inertias.  Each nontrivial projection is represented
naturally by a reflection, and the others are identity.
The natural projective model is faithful on the simple
quotient, by the centralizer/perfectness argument of
D-OL-07.  Hence at most twelve natural reflections
$r_{v_j}$ generate a projective subgroup containing
$S_i/Z(S_i)$.

The space
\[
                         U=\bigcap_j v_j^\perp
                         \quad\text{has dimension at least }d_i-12>0.
                         \tag{OE.22}
\]
Every one of the chosen natural reflections fixes $U$
pointwise.  Consequently every $s\in S_i$ acts on $U$
by a scalar: its projective image is a word in those
reflections.  Those scalars define a homomorphism
$S_i\to F_i^*$, which is trivial since $S_i$ is perfect
\cite[Theorems~11.47--11.48 and proofs, pp.~161--162]{OS:Taylor1992}.
This would make $U$ a nonzero fixed space for the natural
$S_i$-action.  There is no such space: two full opposite
line radicals have natural fixed spaces the two distinct
singular lines, and they generate $S_i$.
This contradiction excludes the source.  If a factor
has $q_i\ne3$, all twelve projections on it are already
trivial, which yields the same contradiction even sooner.

Finally, the selfdual ordinary cases are covered by
D-OL-09; the even-field ordinary cases by D-OL-10--11;
the odd-natural odd-field cases by D-OL-12--15; and the
remaining even-natural odd-field cases by the present
theorem.  Applied in the main transitive setting, these
give the stated ordinary natural-dimension bounds.
None of these implications extends to a spin module
which does not descend to the natural orthogonal group.
\end{proof}
